\section{An auxiliary number process and its finite coagulation correction}
\label{sec:auxiliary}

Fix a solution in \cref{def:solution}. In this section the rates and
daughter kernel retain the full generality of \cref{ass:rates}.
The construction uses the deterministic curve $n_t$ as its environment.
Its one-time laws are the number sampling laws $\eta_t$, but its paths
are not physical mass tags or chronological particle lineages.

\subsection{Construction and identification of the marginals}

\begin{theorem}[Number-weighted jump representation]\label{thm:number-process}
There is a nonexplosive time-inhomogeneous Markov jump process $X_t$
on $E$, with initial law $\eta_0$, whose law at every $t$ is $\eta_t$.
Its generator, meaning its jump kernel acting on bounded Borel
functions $f$, is
\begin{align}
 \mathcal L_t f(x)
 &=\sigma(t)\int_E[f(z)-f(x)]b_{t,x}(\dd z)\notag\\
 &\quad+\lambda(t)N(t)x\int_E[f(x+y)-f(x)]\eta_t(\dd y).
 \label{eq:number-generator}
\end{align}
Thus fragmentation occurs at rate $2\sigma(t)$ and selects a daughter
with law $b_{t,x}/2$. Coagulation occurs at rate $\lambda(t)N(t)x$
and adds a fresh sample from the prescribed law $\eta_t$.
The fragmentation term has a direct interpretation: a fragmentation
event replaces the tagged particle by one of its two daughters chosen
uniformly, which has law $b_{t,x}/2$; number weighting counts both
daughters, so the number-weighted process sees this daughter law at
twice the selection rate $\sigma(t)$.
\end{theorem}

\begin{proof}
Symmetry rewrites the coagulation term of \eqref{eq:weak} as
\[
 \lambda(t)\iint x[f(x+y)-f(x)-f(y)]n_t(\dd x)n_t(\dd y).
\]
For each bounded Borel $f$, divide its scalar, absolutely continuous
balance by $N(t)$ and use \eqref{eq:count}. The result is
\begin{equation}
 \pair{f}{\eta_t}-\pair{f}{\eta_0}
 =\int_0^t\pair{\mathcal L_s f}{\eta_s}\dd s.
 \label{eq:number-forward}
\end{equation}
All these integrals are finite on bounded time intervals.

For the construction use two independent Poisson random measures,
independent also of $X_0$. The fragmentation driver has intensity
$2\sigma(t)\dd t\dd u$, $0<u<1$; the quantile map
$(t,x,u)\mapsto\inf\{z:b_{t,x}((0,z])/2\geq u\}$ chooses its
daughter using $u$. This map is jointly measurable: the set where
it is at most $z$ equals $\{(t,x,u):b_{t,x}((0,z])/2\geq u\}$,
by right-continuity in $z$ of the distribution function, and
$(t,x,z)\mapsto b_{t,x}((0,z])$ is measurable because $b$ is a
measurable kernel. The coagulation driver has
intensity $\dd t\dd v\dd u$, $v>0$, $0<u<1$; retain a point when
$v\leq\lambda(t)N(t)X_{t-}$ and use the quantile map of $\eta_t$,
jointly measurable in $(t,u)$ for the same reason, to choose its
added size. Let $(\mathfrak F_t)_{t\geq0}$ be the filtration
generated by $X_0$ and the two driving Poisson random measures,
augmented with the $\Prob$-null sets; it satisfies the usual
conditions. Because the two drivers are independent of each other
and of $X_0$, each driver keeps its deterministic compensator in
this joint filtration. All stopping times, compensators, and local
martingales in this section refer to $(\mathfrak F_t)$.
Recursive construction
up to successive jumps gives the minimal process, initially defined
up to a possible explosion time. Its jump kernel $Q_t(x,\dd z)$ is
the sum of the two positive kernels in \eqref{eq:number-generator},
and its total rate is
\begin{equation}
 q_t(x)=2\sigma(t)+\lambda(t)N(t)x.
 \label{eq:number-total-rate}
\end{equation}
These rates are bounded by an integrable time function on every
bounded size set and finite time interval.

Here is a first-moment nonexplosion argument. Daughter mass conservation
and $\pair{x}{\eta_t}=m/N(t)$ give, for $V(x)=x$,
\begin{equation}
 \mathcal L_t V(x)=(b(t)-\sigma(t))x\leq b(t)x.
 \label{eq:number-lyapunov}
\end{equation}
The generator was defined on bounded functions; its application to
the unbounded $V$ is justified by the stopping that follows, which
keeps every term of the Dynkin formula finite.
Stop at the $k$th jump $\tau_k$ and at $\tau_R$, the first time
the state exceeds $R$.
Before these stops rates are integrably bounded. An upward jump can
overshoot $R$, but its mark has finite conditional mean $m/N(t)$;
the expected positive and negative increments in the stopped Dynkin
formula are therefore finite. Gronwall's inequality gives
$\E X_{t\wedge\tau_k\wedge\tau_R}\leq\E X_0e^{\Bc(t)}$.
For fixed $k$, only finitely many finite states are visited, so the
size stopping can be removed as $R\to\infty$. Fatou's lemma yields
\begin{equation}
 \E X_{t\wedge\tau_k}\leq\frac m{N_0}e^{\Bc(t)}.
 \label{eq:stopped-first-moment}
\end{equation}
The jump-count compensator is the time integral of the conditional
jump rate; after integrable stopping, its expectation equals the
expected jump count. The compensator of the stopped total jump count has
expectation at most
\begin{equation}
 L_T:=2F(T)+\frac m{N_0}
       \int_0^T\lambda(s)N(s)e^{\Bc(s)}\dd s<\infty.
 \label{eq:stopped-jump-bound}
\end{equation}
With the size stop in place all rates are integrably bounded, so
the compensator identity is exact; as $R\to\infty$ the stopped
counts and their compensators increase to those stopped at
$\tau_k\wedge T$ alone, and monotone convergence in $R$ justifies
the compensator calculation without the size stop. Since
$k\Prob(\tau_k\leq T)\leq L_T$, there is no finite-time explosion.
It also proves nonexplosion from each deterministic finite starting
state. Recursive sampling with a prescribed environment gives the
Markov property.

It remains to identify the marginals; a formal forward equation alone
would not suffice. The prescribed probability curve $\eta_t$ has
integrable jump flux, since
\begin{equation}
 \int_E q_t(x)\eta_t(\dd x)=2\sigma(t)+b(t).
 \label{eq:prescribed-flux}
\end{equation}
Put $G_s(\dd z)=\int_E Q_s(x,\dd z)\eta_s(\dd x)$ and
$S_{s,t}(z)=\exp(-\int_s^t q_u(z)\dd u)$. The weak loss identity
\eqref{eq:loss-duhamel} of \cref{app:existence}, used with weight
one, turns \eqref{eq:number-forward} into
\begin{equation}
 \eta_t(\dd z)=S_{0,t}(z)\eta_0(\dd z)
       +\int_0^t S_{s,t}(z)G_s(\dd z)\dd s.
 \label{eq:number-survival}
\end{equation}
What is used is the following. Restrict the output variable to
$z\in(0,R]$; there the loss multiplier $q_s(z)$ is bounded by the
integrable function $2\sigma(s)+\lambda(s)N(s)R$, and the incoming
gain is the restriction of the \emph{full} $G_s$, including arrivals
from parents larger than $R$. The restricted loss equation has a
unique solution with bounded total variation on bounded time
intervals, given by the exponential formula: existence is the
convergent loss-operator series, and uniqueness is Gronwall's
inequality for the variation of the difference of two solutions.
This proves the formula on the output set $(0,R]$, and positivity
allows $R\to\infty$. Both gain and loss have finite integrated
total mass by \eqref{eq:prescribed-flux}. This argument uses weak
kernel integrals, not a Bochner derivative.

The right-hand side of \eqref{eq:number-survival} is positive.
Its successive substitutions show that $\eta_t$ dominates the sum of
the zero-jump, one-jump, and all subsequent finite-jump contributions
of the minimal construction. More explicitly, start with
$\nu_t^{(0)}=S_{0,t}\eta_0$ and define
\[
 \nu_t^{(j+1)}(\dd z)=\int_0^t S_{s,t}(z)
                 (Q_s^*\nu_s^{(j)})(\dd z)\dd s.
\]
Then $\eta_t\geq\sum_{j=0}^k\nu_t^{(j)}$ for every $k$.
Nonexplosion makes $\sum_{j\geq0}\nu_t^{(j)}$ a probability measure,
the law of $X_t$. Domination between probability measures forces
equality. No additional uniqueness theorem for an unbounded forward
operator, or second size moment, is needed.
\end{proof}

\begin{remark}[Relation to mass tagging and to classical constructions]
\label{rem:inverse-size}
None of the ingredients of \cref{thm:number-process} is new, and we
state plainly what is classical. The fragmentation part, rate
$2\sigma$ with daughter law $b_{t,x}/2$, is the classical
many-to-one or intensity-measure computation for finite-activity
conservative binary fragmentation: the expected number of particles
in a Borel set evolves linearly, and its normalization by the total
count is the law of a tag that follows a uniformly chosen daughter.
The coagulation part is the number-weighted counterpart of the mass
tagging of \citet{DeaconuFournierTanre2002}: the number-weighted
process is the $h=1/x$ Doob transform of that mass-weighted process,
as the identity below shows. For pure coagulation the auxiliary path
is the cluster growth chain associated with the additive coalescent
of \citet{AldousPitman1998}, seen through the total progeny of
Poisson Galton--Watson trees (cf.\ \citet{Bertoin2002} and
\cref{cor:pure-coag}). The new ingredient of this section is only
the finite correction bound of \cref{thm:controlled-path}.

Jump representations of mass sampling in coagulation are established;
see \citet{DeaconuFournierTanre2002}. The corresponding mass-tag
generator, with the present fragmentation added, is
\[
 \mathcal L_t^{\rm mass}f(x)
 =\lambda(t)\int_E(x+y)[f(x+y)-f(x)]n_t(\dd y)
   +\frac{\sigma(t)}x\int_E z[f(z)-f(x)]b_{t,x}(\dd z).
\]
For $h(x)=1/x$, direct integration gives
$\mathcal L_t^{\rm mass}h=(\sigma(t)-b(t))h$ and the algebraic
inverse-size Doob-transform identity
\[
 h^{-1}\mathcal L_t^{\rm mass}(hf)
       -(\sigma(t)-b(t))f=\mathcal L_t f.
\]
This calculation explains the change of sampling weight and the
fragmentation rate $2\sigma$. It is an identity of generators;
the preceding direct construction supplies the required path law.
In particular the mass-tag coagulation rate is
$\lambda(t)N(t)x+b(t)$, whereas the number process has rate
$\lambda(t)N(t)x$. Conclusions about the number process below are
specific to this representation.
\end{remark}

\subsection{A finite logarithmic correction under arbitrary controls}

Let $s_j$ be the fragmentation times and
$\Theta_j=X_{s_j}/X_{s_j-}\in(0,1)$ their retained fractions.
For a coagulation time $s$, write $U_s=X_s-X_{s-}>0$ and set
\begin{equation}
 A_t=\sum_{\substack{s\leq t\\s\text{ coagulation}}}
             \log(1+U_s/X_{s-}),\qquad
 a(t)=\frac{\lambda(t)}{N(t)}\iint_{E^2}
           x\log(1+y/x)n_t(\dd x)n_t(\dd y).
 \label{eq:log-correction}
\end{equation}

\begin{theorem}[Finite correction and finite total coagulations]
\label{thm:controlled-path}
For the process in \cref{thm:number-process},
\begin{equation}
 \log X_t=\log X_0+\sum_{s_j\leq t}\log\Theta_j+A_t.
 \label{eq:path-log-identity}
\end{equation}
The nonnegative process $A_t$ increases to a finite $A_\infty$ almost
surely and in $L^1$, with $\E A_t=\int_0^t a(s)\dd s$ and
\begin{align}
 \E(A_\infty-A_t)
 &\leq\frac{H(0)^2}{mN_0}
                  \int_t^\infty b(s)e^{-\kappa C(s)}\dd s\notag\\
 &\leq\frac{H(0)^2}{\kappa mN_0}e^{-\kappa C(t)}.
 \label{eq:controlled-log-tail}
\end{align}
In particular $\E A_\infty\leq1/\kappa$.
The total number $K_\infty$ of auxiliary coagulation jumps is finite
almost surely, although its mean is exactly
\begin{equation}
 \E K_t=\Bc(t),\qquad \E K_\infty=\Bc(\infty),
 \label{eq:coag-count-mean}
\end{equation}
which is infinite whenever $\Bc(\infty)=\infty$, in particular for
every constant $\lambda>0$. No initial or daughter logarithmic moment
is assumed.
\end{theorem}

\begin{proof}
Nonexplosion makes \eqref{eq:path-log-identity} an identity of finite
sums at each finite time. The elementary inequality
$\log(1+r)\leq\sqrt r$ for $r>0$ follows by differentiating their
difference: the derivative is
$-(\sqrt r-1)^2/[2\sqrt r(1+r)]\leq0$, and the difference tends
to zero at zero. Thus \cref{thm:affinity} gives
\begin{equation}
 0\leq a(t)\leq\frac{\lambda(t)H(t)^2}{N(t)}
 \leq\frac{H(0)^2}{mN_0}b(t)e^{-\kappa C(t)}.
 \label{eq:log-rate-bound}
\end{equation}
The nonnegative jump compensator and marginal identification give
$\E A_t=\int_0^t a(s)\dd s$. Since $C'=b+\sigma\geq b$,
\[
 \int_t^\infty b(s)e^{-\kappa C(s)}\dd s
 \leq\frac{e^{-\kappa C(t)}-e^{-\kappa C(\infty)}}{\kappa}
 \leq\frac1\kappa e^{-\kappa C(t)},
\]
where $e^{-\infty}=0$. Monotone convergence proves finiteness of
$\E A_\infty$ and \eqref{eq:controlled-log-tail}; it also gives
$\E|A_\infty-A_t|\to0$. When $C(\infty)<\infty$, the final
exponential bound need not tend to zero, but the preceding integral
bound does. Finally $H(0)^2\leq mN_0$.

To prove finite total coagulation, define
\begin{equation}
 \mathcal M_t=e^{F(t)}\prod_{s_j\leq t}\Theta_j.
 \label{eq:daughter-martingale}
\end{equation}
Conditional on the full past and on a fragmentation at $(t,x)$,
the multiplier has mean
\[
 \frac1{2x}\int_E z b_{t,x}(\dd z)=\frac12.
\]
The predictable drift of \eqref{eq:daughter-martingale} is therefore
$\sigma(t)\mathcal M_{t-}+2\sigma(t)\mathcal M_{t-}(1/2-1)=0$.
It is a local martingale in the construction filtration
$(\mathfrak F_t)$, which includes the coagulation randomness. Since
$0<\mathcal M_t\leq e^{F(T)}$ on $[0,T]$, it is a true mean-one
martingale on every finite horizon. Doob's maximal inequality gives
$\Prob(\sup_{t\leq T}\mathcal M_t>L)\leq1/L$; first let
$T\to\infty$ and then $L\to\infty$ to obtain
$\sup_{t\geq0}\mathcal M_t<\infty$ almost surely.

Combining \cref{eq:count,eq:path-log-identity,eq:daughter-martingale},
\[
 N(t)X_t=N_0X_0e^{-\Bc(t)}\mathcal M_t e^{A_t}.
\]
Hence the cumulative predictable coagulation intensity
$\Lambda_t:=\int_0^t\lambda(s)N(s)X_{s-}\dd s$ satisfies
\begin{align}
 \Lambda_\infty&:=\lim_{t\to\infty}\Lambda_t\notag\\
 &\leq\frac{N_0X_0}{m}e^{A_\infty}
       \sup_{t\geq0}\mathcal M_t
       \int_0^\infty b(t)e^{-\Bc(t)}\dd t<\infty
       \quad\text{almost surely},
 \label{eq:finite-coag-hazard}
\end{align}
because the last integral is $1-e^{-\Bc(\infty)}\leq1$.
All other factors are finite pathwise; their expectations are not needed.
For completeness, stop the coagulation count at $\tau_L$, the first
time its continuous compensator $\Lambda_t$ reaches a positive level
$L$; this is an $(\mathfrak F_t)$-stopping time. The expected stopped
count over all time is at most $L$, by compensation and monotone
convergence, so it is finite almost surely. On
$\{\Lambda_\infty<L\}$ this stop changes no coagulation jump.
The union over integer $L$ proves $K_\infty<\infty$ almost surely.
On the other hand, marginal identification and Tonelli give
\[
 \E K_t=\int_0^t\lambda(s)N(s)\E X_{s-}\dd s
       =\int_0^t b(s)\dd s=\Bc(t).
\]
There are no fixed-time jumps under the absolutely continuous
driving intensities, so the left limits do not change this integral.
Monotone convergence as $t\to\infty$ gives $\E K_\infty=\Bc(\infty)$,
proving \eqref{eq:coag-count-mean}; for constant $\lambda>0$ this
is $\E K_\infty=\infty$.
\end{proof}

After its finite last coagulation, $A_t$ is exactly constant. That last
time generally is not a stopping time, so this assertion does not make
the subsequent fragmentation independent of it. Nor does finite
$K_\infty$ imply that coagulation ceases in the population. Since
$\E K_\infty=\Bc(\infty)$, the mean total count is infinite whenever
$\Bc(\infty)=\infty$, in particular for every constant $\lambda>0$:
rare auxiliary paths account for an infinite mean even though each
path has only finitely many coagulation events.

\begin{corollary}[A pathwise dichotomy for the constructed auxiliary
process]
\label{cor:controlled-size}
For the auxiliary process constructed in \cref{thm:number-process}:
if $F(t)\to\infty$, then $X_t\to0$ almost surely and
$\eta_t\Rightarrow\delta_0$ on $[0,\infty)$; if $F(\infty)<\infty$,
then $X_t$ is eventually constant at a positive finite random value,
and $\eta_t$ converges weakly on $E$ to its law.
The equation determines only the one-time laws $\eta_t$, so the
weak limits of $\eta_t$ are statements about the equation, whereas
the almost-sure path statements concern the multi-time law of the
constructed process, which is an artifact of the construction.
\end{corollary}

\begin{proof}
The path product gives
$X_t=X_0e^{-F(t)}\mathcal M_t e^{A_t}$, with the last two factors
bounded almost surely. This proves the first assertion. In the second
case the fragmentation driver has finite total intensity $2F(\infty)$,
so it has finitely many points almost surely. Together with
$K_\infty<\infty$ this leaves only finitely many jumps of either type,
each between positive finite sizes. The weak limits follow by bounded
convergence applied to continuous tests.
\end{proof}
